%\RequirePackage{kvoptions-patch}
\documentclass[a4paper, twoside, bibliography=totoc, headsepline, cleardoublepage=empty, parskip=half, draft=false]{scrbook}

\let\ifdeutsch\iffalse
\let\ifenglish\iftrue

\input{config}


\usepackage[
  title={This is a Two Line Long Title it is Filling Two Lines}, % Do not forget to capitalize your title correctly, you may use the following page to help you: https://capitalizemytitle.com/
  author={Sigmund Freud},
  orcid=0000-0000-0000-0000, % get your own ORCID via https://orcid.org/
  email={sigmund.freud@ifi.lmu.de},
  type=bachelor,
  institute={Department of Computer Science}, % Institut für Informatik
  course={Computer Science},
  examinerFirst={Prof.\ Dr.\ Uwe Fessor},
  examinerSecond={Prof.\ Dr.\ Uwe Fessor},
  supervisor={}, %Dipl.-Inf.\ Roman Tiker,\\Dipl.-Inf.\ Laura Stern,\\Otto Normalverbraucher,\ M.Sc.
  startdate={July 5, 2018},
  enddate={January 5, 2019},
  copyright=ccbysa, % ccbysa, ccbynosa, cc0, nda, none
  language=english,
  aids=true, % Specify if you used any aids in the crated of the thesis, if yes [true] if no aids were used [false]
  % Next, enter which tools were used and how according to the following taxonomy (https://ceur-ws.org/GenAI/Taxonomy.html). If no aid was used, the next three values can be ignored.
  aidstools={Gammerly, ChatGPT}, 
  aidscontributions={Grammar and spelling check, Paraphrase and reword},
  aidsextracontributions={Dall-E 2 was used to generate Figure X. }
]{tudo-thesis-cover}

\input{acronyms}

\makenoidxglossaries

\begin{document}

\input{commands}

\pagenumbering{arabic}
\Coverpage
\Copyright
%Eigener Seitenstil fuer die Kurzfassung und das Inhaltsverzeichnis
\deftriplepagestyle{preamble}{}{}{}{}{}{\pagemark}
%Doku zu deftriplepagestyle: scrguide.pdf
\pagestyle{preamble}
\renewcommand*{\chapterpagestyle}{preamble}



%Kurzfassung / abstract
%auch im Stil vom Inhaltsverzeichnis
\section*{Kurzfassung}

\todo{Short summary of the thesis. Here, the following questions should be answered:}
\todo{What is the specific problem addressed?}
\todo{What have you done?}
\todo{What did you find out?}
\todo{What are the implications on a larger scale?}
\todo{Should be around 0.5 pages. Not longer than 1 page.}

\cleardoublepage

\section*{Abstract}

\todo{Short summary of the thesis. Here, the following questions should be answered:}
\todo{What is the specific problem addressed?}
\todo{What have you done?}
\todo{What did you find out?}
\todo{What are the implications on a larger scale?}
\todo{Should be around 0.5 pages. Not longer than 1 page.}

\cleardoublepage


% BEGIN: Verzeichnisse

\iftex4ht
\else
  \microtypesetup{protrusion=false}
\fi

%%%
% Literaturverzeichnis ins TOC mit aufnehmen, aber nur wenn nichts anderes mehr hilft!
% \addcontentsline{toc}{chapter}{Literaturverzeichnis}
%
% oder zB
%\addcontentsline{toc}{section}{Abkürzungsverzeichnis}
%
%%%

%Produce table of contents
%
%In case you have trouble with headings reaching into the page numbers, enable the following three lines.
%Hint by http://golatex.de/inhaltsverzeichnis-schreibt-ueber-rand-t3106.html
%
%\makeatletter
%\renewcommand{\@pnumwidth}{2em}
%\makeatother
%
\tableofcontents

% Bei einem ungünstigen Seitenumbruch im Inhaltsverzeichnis, kann dieser mit
% \addtocontents{toc}{\protect\newpage}
% an der passenden Stelle im Fließtext erzwungen werden.

\listoffigures
\listoftables

% Control List of Listings
\let\iflistings\iffalse
%Wird nur bei Verwendung von der lstlisting-Umgebung mit dem "caption"-Parameter benoetigt
%\lstlistoflistings
%ansonsten:
\iflistings
  \ifdeutsch
    \listof{Listing}{Verzeichnis der Listings}
  \else
    \listof{Listing}{List of Listings}
  \fi
\fi

% Control List of Algorithms
\let\ifalgorithms\iffalse
\ifalgorithms
  %mittels \newfloat wurde die Algorithmus-Gleitumgebung definiert.
  %Mit folgendem Befehl werden alle floats dieses Typs ausgegeben
  \ifdeutsch
    \listof{Algorithmus}{Verzeichnis der Algorithmen}
  \else
    \listof{Algorithmus}{List of Algorithms}
  \fi
  %\listofalgorithms %Ist nur für Algorithmen, die mittels \begin{algorithm} umschlossen werden, nötig
\fi

% Control Glossary
\let\ifglossary\iftrue
\ifglossary
    \ifdeutsch
        \printnoidxglossary[type=\acronymtype, title=Abkürzungsverzeichnis]
    \else
        \printnoidxglossary[type=\acronymtype, title=List of Acronyms]
    \fi
\fi

\iftex4ht
\else
  %Optischen Randausgleich und Grauwertkorrektur wieder aktivieren
  \microtypesetup{protrusion=true}
\fi
% END: Verzeichnisse


% Headline and footline
\renewcommand*{\chapterpagestyle}{scrplain}
\pagestyle{scrheadings}
\pagestyle{scrheadings}
\ihead[]{}
\chead[]{}
\ohead[]{\headmark}
\cfoot[]{}
\ofoot[\usekomafont{pagenumber}\thepage]{\usekomafont{pagenumber}\thepage}
\ifoot[]{}


%% vv  scroll down for content  vv %%
\input{content.tex}


%TC:ignore
\printbibliography

All links were last followed on \today{}.


\appendix
\input{latexhints/latexhints-english}
% How to make use of an Appendix?
% The appendix is the place for material that is important for transparency, verification, and reproducibility, but not necessary for following the main argument of the thesis. The main chapters should tell a clear and readable scientific story, explaining the problem, the method, and the results. The appendix exists to hold all technical, methodological, or detailed material that would interrupt this flow if placed in the main text.
% Typical content of an appendix includes extended tables and figures, full questionnaires or interview guides, detailed statistical analyses, additional results, algorithm descriptions, model architectures, hyperparameters, and implementation details. A reader should be able to understand and evaluate the thesis without reading the appendix, but a technically interested reader should be able to reproduce and critically inspect the work using the information provided there.
% You should refer to the appendix from the main text whenever relevant, for example “see Appendix A for details of the model architecture” or “additional results are provided in Appendix A”. The appendix is not a place to hide weak or unimportant results. It exists to strengthen the credibility and scientific rigor of your work by making all relevant details openly available.

\Aids

\Affirmation
%TC:endignore
\end{document}
